\begin{flushleft}
{\Large
\textbf{DeepStarNoodles: A deep-learning pipeline for high-throughput behavioral analysis of sunflower stars}
}
\newline
\\
Willem Weertman\textsuperscript{1,2*},
Dalton Blackwell\textsuperscript{2},
Chloe Schwab\textsuperscript{2},
Julia Knopf\textsuperscript{2},
Tess Chapman\textsuperscript{2},
Pema Kitaeff\textsuperscript{2},
Brooke Ashcraft\textsuperscript{2},
Rebecca Guenther\textsuperscript{2},
Jason Hodin\textsuperscript{2}
\\
\bigskip
\textbf{1} Department of Psychology, University of Washington, Seattle, Washington, United States of America
\\
\textbf{2} Friday Harbor Laboratories, University of Washington, Friday Harbor, Washington, United States of America
\\
\bigskip

* willemlw@uw.edu

\end{flushleft}

\section*{Abstract}

Pose-estimation tools locate a fixed list of body points, which does not describe an animal whose number of arms varies.
The sunflower sea star, \textit{Pycnopodia helianthoides}, is such an animal: individuals in our images have 3 to 28 arms that look alike and bend along their whole length.
The species is Critically Endangered and is now reared in captivity, where behaviour is scored by observers.
We describe DeepStarNoodles, a pipeline that returns one base-to-tip polyline for each arm of each animal in a frame and links arms from frame to frame.
A detector outlines each animal, a second network predicts 18 dense maps of arm geometry on the crop, and a recurrent network with 194,463 parameters walks along each arm from tip to base.
Training used 7,978 images with instance masks and 5,093 crops with 57,055 annotated arms.
On test images the detector had precision 0.920 and recall 0.903; recall was 0.924 on laboratory images and 0.668 on field images apart from one cruise that also supplied training images.
The tracer reached an F1 of 0.755 averaged over 180 acquisition groups, against 0.705 for a shortest-path decoder of the same maps (difference 5.0 percentage points, 95\% CI 4.0 to 6.0).
Matched arms lay a median of 2.7~px from the annotation in a 512~px crop.
Accuracy was lowest for animals with few arms.
On one rented GPU the pipeline processed 58 frames per second of streamed video.
In recordings of animals righting themselves after inversion, righting times read from the pipeline outputs had a mean absolute difference of 2.4~s from an observer's over 273 bouts.
Across 96 changes from a quiescent to an alert state, arm-tip motion rose by a median of 63\%.
Code is at \url{https://github.com/weertman/DeepStarNoodles}; models and annotated data are on Hugging Face.

\section*{Author summary}

Sunflower sea stars were once common along the Pacific coast of North America, and a disease outbreak that began in 2013 killed most of them.
Laboratories now raise the species, and we film the animals to learn how they behave.
Turning that video into measurements is the slow step.
Software built for mice or flies follows a fixed set of body points.
A sunflower sea star has no fixed set: one animal may have eight arms and another more than twenty, and every arm can bend anywhere along its length.
We built software that finds each animal in a video frame and draws a line along each of its arms.
It first marks where arms are and which way they run, and then walks along each arm in small steps, as a person might trace it with a pencil.
We measured how often the traced arms agree with arms drawn by people, where the software fails, and how fast it runs.
As examples, we timed how long overturned animals took to right themselves, and measured how arm movement changed as resting animals became alert.
We share the software, the trained models and the annotated images.
